% =====================================================================
% Chapter 5 — Logic and Puzzles — part a: theory, examples, problems 1–13
% =====================================================================
% helper macros (prefix \lga..., guarded by \providecommand)
% \lgaScale{left}{right}: a balanced two-pan scale
\providecommand{\lgaScale}[2]{\begin{tikzpicture}[baseline=-0.3cm,font=\small]
  \draw[Ink,line width=1.1pt] (-2.25,0)--(-0.45,0) (0.45,0)--(2.25,0);
  \draw[Ink,line width=1pt] (-1.35,0)--(-1.35,-0.5)--(1.35,-0.5)--(1.35,0);
  \draw[Ink,line width=1pt] (-0.17,-0.78)--(0,-0.5)--(0.17,-0.78);
  \node[above=0.5pt,inner sep=1pt] at (-1.35,0){#1};\node[above=0.5pt,inner sep=1pt] at (1.35,0){#2};\end{tikzpicture}}
% shapes used on the scales
\providecommand{\lgaSq}{\tikz[baseline=0pt]\draw[fill=tmBlue!65,draw=tmBlue!55!black,line width=0.6pt] (0,0) rectangle (0.3,0.3);}
\providecommand{\lgaTri}{\tikz[baseline=0pt]\draw[fill=tmOrange!85,draw=tmOrange!55!black,line width=0.6pt] (0,0)--(0.34,0)--(0.17,0.31)--cycle;}
\providecommand{\lgaCir}{\tikz[baseline=0pt]\draw[fill=tmGreen!65,draw=tmGreen!45!black,line width=0.6pt] (0.15,0.15) circle (0.15);}
% coloured boxes (Problem 23)
\providecommand{\lgaBox}[1]{\tikz[baseline=0pt]\draw[fill=#1!70,draw=#1!55!black,line width=0.6pt,rounded corners=1pt] (0,0) rectangle (0.34,0.3);}
% \lgaMagic{a}{b}{c}{d}{e}{f}{g}{h}{i}: 3x3 grid (row by row)
\providecommand{\lgaMagic}[9]{\begin{tikzpicture}[baseline=(current bounding box.center),font=\normalsize]
  \fill[tmBlue!5] (0,0) rectangle (2.7,2.7);
  \draw[Navy,line width=0.9pt] (0,0) grid[step=0.9] (2.7,2.7);
  \node at (0.45,2.25){#1};\node at (1.35,2.25){#2};\node at (2.25,2.25){#3};
  \node at (0.45,1.35){#4};\node at (1.35,1.35){#5};\node at (2.25,1.35){#6};
  \node at (0.45,0.45){#7};\node at (1.35,0.45){#8};\node at (2.25,0.45){#9};\end{tikzpicture}}
% \lgaFour{c/r/value,...}: 4x4 grid split into four 2x2 boxes (c = column 1..4 from the left, r = row 1..4 from the top)
\providecommand{\lgaFour}[1]{\begin{tikzpicture}[baseline=(current bounding box.center),font=\normalsize]
  \fill[tmBlue!4] (0,0) rectangle (2.8,2.8);
  \draw[Navy!60,line width=0.5pt] (0,0) grid[step=0.7] (2.8,2.8);
  \draw[Navy,line width=1.6pt] (0,0) rectangle (2.8,2.8) (1.4,0)--(1.4,2.8) (0,1.4)--(2.8,1.4);
  \foreach \c/\r/\v in {#1}{\node at ({0.7*\c-0.35},{2.8-0.7*\r+0.35}){\v};}\end{tikzpicture}}
% \lgaPips{n}: the pips of a die face in the unit square
\providecommand{\lgaPips}[1]{\ifcase#1\or
  \fill (0.5,0.5) circle (0.085);\or
  \fill (0.27,0.73) circle (0.085) (0.73,0.27) circle (0.085);\or
  \fill (0.25,0.75) circle (0.085) (0.5,0.5) circle (0.085) (0.75,0.25) circle (0.085);\or
  \fill (0.27,0.27) circle (0.085) (0.73,0.73) circle (0.085) (0.27,0.73) circle (0.085) (0.73,0.27) circle (0.085);\or
  \fill (0.25,0.25) circle (0.085) (0.75,0.75) circle (0.085) (0.25,0.75) circle (0.085) (0.75,0.25) circle (0.085) (0.5,0.5) circle (0.085);\or
  \fill (0.27,0.22) circle (0.08) (0.73,0.22) circle (0.08) (0.27,0.5) circle (0.08) (0.73,0.5) circle (0.08) (0.27,0.78) circle (0.08) (0.73,0.78) circle (0.08);\fi}
% \lgaCube{top}{front}{right}: a die seen at an angle (3 faces visible)
\providecommand{\lgaCube}[3]{\begin{tikzpicture}[baseline=0.4cm,scale=0.85,line join=round]
  \draw[Ink,line width=0.9pt,fill=white] (0,0) rectangle (1,1);
  \draw[Ink,line width=0.9pt,fill=white] (0,1)--(1,1)--(1.45,1.45)--(0.45,1.45)--cycle;
  \draw[Ink,line width=0.9pt,fill=white] (1,0)--(1.45,0.45)--(1.45,1.45)--(1,1)--cycle;
  \begin{scope}[Ink]\lgaPips{#2}\end{scope}
  \begin{scope}[Ink,cm={1,0,0.45,0.45,(0,1)}]\lgaPips{#1}\end{scope}
  \begin{scope}[Ink,cm={0.45,0.45,0,1,(1,0)}]\lgaPips{#3}\end{scope}\end{tikzpicture}}
% \lgaTriEight{apex}{l1}{l2}{bottom left}{r1}{r2}{bottom right}{bottom middle}: triangle with 8 circles
\providecommand{\lgaTriEight}[8]{\begin{tikzpicture}[baseline=(current bounding box.center),font=\small]
  \draw[Navy,line width=0.9pt] (0,2.7)--(-1.5,0)--(1.5,0)--cycle;
  \foreach \x/\y in {0/2.7,-0.5/1.8,-1/0.9,-1.5/0,0.5/1.8,1/0.9,1.5/0,0/0}{\draw[Navy,fill=white,line width=0.9pt] (\x,\y) circle (0.27);}
  \node at (0,2.7){#1};\node at (-0.5,1.8){#2};\node at (-1,0.9){#3};\node at (-1.5,0){#4};
  \node at (0.5,1.8){#5};\node at (1,0.9){#6};\node at (1.5,0){#7};\node at (0,0){#8};\end{tikzpicture}}
% \lgaTriSix{apex}{bottom left}{bottom right}{middle of left side}{middle of base}{middle of right side}: triangle with 6 circles
\providecommand{\lgaTriSix}[6]{\begin{tikzpicture}[baseline=(current bounding box.center),font=\small]
  \draw[Navy,line width=0.9pt] (0,2.2)--(-1.27,0)--(1.27,0)--cycle;
  \foreach \x/\y in {0/2.2,-1.27/0,1.27/0,-0.635/1.1,0/0,0.635/1.1}{\draw[Navy,fill=white,line width=0.9pt] (\x,\y) circle (0.27);}
  \node at (0,2.2){#1};\node at (-1.27,0){#2};\node at (1.27,0){#3};
  \node at (-0.635,1.1){#4};\node at (0,0){#5};\node at (0.635,1.1){#6};\end{tikzpicture}}
% \lgaSqEight{TL}{top}{TR}{right}{BR}{bottom}{BL}{left}: square with 8 circles (clockwise)
\providecommand{\lgaSqEight}[8]{\begin{tikzpicture}[baseline=(current bounding box.center),font=\small]
  \draw[Navy,line width=0.9pt] (0,0) rectangle (2.4,2.4);
  \foreach \x/\y in {0/2.4,1.2/2.4,2.4/2.4,2.4/1.2,2.4/0,1.2/0,0/0,0/1.2}{\draw[Navy,fill=white,line width=0.9pt] (\x,\y) circle (0.27);}
  \node at (0,2.4){#1};\node at (1.2,2.4){#2};\node at (2.4,2.4){#3};\node at (2.4,1.2){#4};
  \node at (2.4,0){#5};\node at (1.2,0){#6};\node at (0,0){#7};\node at (0,1.2){#8};\end{tikzpicture}}
% \lgaCross{top}{left}{centre}{right}{bottom}: plus-shaped figure of 5 cells
\providecommand{\lgaCross}[5]{\begin{tikzpicture}[baseline=(current bounding box.center),font=\normalsize]
  \foreach \x/\y in {1/2,0/1,1/1,2/1,1/0}{\draw[Navy,fill=tmBlue!5,line width=0.9pt] (\x*0.8,\y*0.8) rectangle ++(0.8,0.8);}
  \node at (1.2,2.0){#1};\node at (0.4,1.2){#2};\node at (1.2,1.2){#3};\node at (2.0,1.2){#4};\node at (1.2,0.4){#5};\end{tikzpicture}}
% \lgaWheel{k}: the 26 letters clockwise around a circle, arrow showing a shift of k from A
\providecommand{\lgaWheel}[1]{\begin{tikzpicture}[baseline=(current bounding box.center),font=\scriptsize\sffamily]
  \draw[Navy!50,line width=0.8pt] (0,0) circle (1.95);
  \foreach \l [count=\i from 0] in {A,B,C,D,E,F,G,H,I,J,K,L,M,N,O,P,Q,R,S,T,U,V,W,X,Y,Z}{
    \node[circle,fill=white,draw=Navy!60,inner sep=0.6pt,minimum size=0.36cm] at ({90-\i*360/26}:1.95){\l};}
  \draw[->,tmOrange,line width=1pt] ({90-3}:1.4) arc ({90-3}:{90-#1*360/26+3}:1.4);
  \node[font=\footnotesize\bfseries,text=tmOrange!80!black] at (0,0){$+#1$};\end{tikzpicture}}


\begin{knowbox}
\textbf{Big idea.} Logical thinking is \emph{the art of drawing certain conclusions from what we already know}. A logic puzzle
usually needs no hard calculation: the difficulty lies in \emph{reading carefully}, \emph{recording information systematically} and
\emph{not missing any case}. In this chapter you will get to know the following ``tools'':
\begin{itemize}[leftmargin=1.4em,itemsep=1pt,topsep=2pt]
  \item \textbf{truth-tellers and liars}, ``half-true'' statements, and how to negate a statement;
  \item \textbf{case analysis} and \textbf{proof by contradiction} (assume, then look for a contradiction);
  \item \textbf{ordering} with arrow diagrams and number lines; telling ``\emph{must be true}'' apart from ``\emph{could be true}'';
  \item \textbf{shift ciphers} (the Caesar cipher) and \textbf{remainders} on division by $26$ (or by $7$ for the days of the week);
  \item \textbf{magic squares}, \textbf{balance scales} (which are really equations), \textbf{sums on figures} and \textbf{dice};
  \item \textbf{Latin squares} and $4\times 4$ \textbf{sudoku}.
\end{itemize}
Four golden habits: \emph{(1) write it down -- use symbols; (2) draw it -- diagrams, tables; (3) consider every case; (4) check the
answer against \textbf{all} the conditions of the problem.}
\end{knowbox}
\vspace{0.3cm}

% ---------------------------------------------------------------------
\needspace{9\baselineskip}\shead{True and false statements: truth-tellers and liars}
\begin{knowbox}
\textbf{1. Statements.} A sentence that is either \emph{true} or \emph{false} (it cannot be both) is called a statement.
Examples: ``$7$ is odd'' (true), ``Paris is the capital of Italy'' (false).

\textbf{2. Truth-tellers and liars.} In puzzles, everything a \emph{truth-teller} says is true, and everything a \emph{liar} says is
false. Therefore:
\begin{itemize}[leftmargin=1.4em,itemsep=1pt,topsep=2pt]
  \item if we know who the speaker is, we know whether the statement is true or false;
  \item conversely, if we know whether the statement is true or false, we know who the speaker is.
\end{itemize}

\textbf{3. Negating a statement (saying ``the opposite'').}
\begin{center}
\renewcommand{\arraystretch}{1.25}
\begin{tabular}{|l|l|}
\hline
\rowcolor{tmBlue!10} Statement & Negation (true exactly when the statement is false) \\ \hline
Anna \emph{and} Ben are both liars & \emph{At least one} of the two tells the truth \\ \hline
\emph{Every} apple is red & \emph{There is at least one} apple that is not red \\ \hline
The flower is a red rose & The flower is not red, \emph{or} it is not a rose (or both) \\ \hline
\end{tabular}
\end{center}

\textbf{4. ``Half-true'' statements.} If a statement has two parts and we know that \emph{exactly one} part is true, then:
first part true $\Rightarrow$ second part false; \ first part false $\Rightarrow$ second part true. Decide true/false for \emph{one}
part, and all the other parts follow like a row of dominoes.

\textbf{5. Lying on certain days.} Someone who ``lies on Tuesdays'' says only false things on Tuesday. To find on which days they
\emph{could} make a given statement, make a table: for each day, decide whether the statement is true or false, and compare with
whether the person tells the truth or lies on that day.
\end{knowbox}
\begin{proofbox}
Why can nobody say ``I am a liar''? If the speaker were a truth-teller, the statement would be true, so they would be a liar --
impossible. If the speaker were a liar, the statement would be false, so they would \emph{not} be a liar -- impossible too. Both
cases lead to a contradiction, so the sentence can never be spoken (on an ``island'' with only truth-tellers and liars).
\end{proofbox}
\begin{warnbox}
The negation of ``Anna and Ben are both liars'' is \emph{not} ``Anna and Ben both tell the truth''. As soon as \emph{one} of them tells
the truth, the original statement is already false.
\end{warnbox}
\begin{tipbox}
For a truth-and-lies puzzle, draw a small table: each row is one possibility (who tells the truth, who lies), each column is one
statement. Fill each cell with T/F; the row with no contradiction is the answer.
\end{tipbox}

% ---------------------------------------------------------------------
\needspace{9\baselineskip}\shead{Case analysis and proof by contradiction}
\begin{knowbox}
\textbf{1. Case analysis.} When something is unknown, we split into \emph{all} the possibilities (missing none) and examine them
one by one. A possibility that leads to something impossible (a contradiction with the problem) is ruled out. The possibility that
remains is the answer.

\textbf{2. Proof by contradiction.} To show that a statement $P$ is true, we \emph{assume} that $P$ is false and keep reasoning.
If we arrive at something impossible, the assumption was wrong, so $P$ is true.

\textbf{3. Two contradictory statements.} If two statements cannot both be true and cannot both be false (for example Leo says
``Helen is lying'' about Helen's statement), then \emph{exactly one} of them is true. This is the ``key'' to very many puzzles: we
immediately know that one statement of the pair is true and the other false.

\textbf{4. Seating around a round table.} Rotating the whole table does not change who sits next to whom. So we may \emph{fix} the
seat of one person and only consider the seats of the others. Around a round table with $5$ seats, each person has exactly $2$
neighbours, and exactly $2$ people who are \emph{not} next to them.
\end{knowbox}
\begin{figC}
  \foreach \i in {0,...,4}{\coordinate (P\i) at ({90-\i*72}:1.35);}
  \draw[Navy!40,line width=5pt] (0,0) circle (0.75);
  \foreach \i/\n in {0/1,1/2,2/3,3/4,4/5}{\node[circle,draw=Navy,fill=tmBlue!8,line width=0.8pt,minimum size=0.55cm] at (P\i){\n};}
  \draw[tmGreen!70!black,line width=1pt] (P0) to[bend left=25] (P1);
  \draw[tmGreen!70!black,line width=1pt] (P0) to[bend right=25] (P4);
  \draw[AnsRed!70,dashed,line width=1pt] (P0)--(P2) (P0)--(P3);
  \node[align=left,font=\footnotesize,text=Slate,anchor=west] at (1.6,0.3){\textcolor{tmGreen!60!black}{solid}: seat $1$ is next to seats $2$ and $5$};
  \node[align=left,font=\footnotesize,text=Slate,anchor=west] at (1.6,-0.3){\textcolor{AnsRed!80}{dashed}: seat $1$ is not next to seats $3$, $4$};
\end{figC}
\begin{proofbox}
Why is proof by contradiction valid reasoning? A statement can only be true or false. If ``$P$ is false'' leads to something
impossible, then the possibility ``$P$ is false'' cannot happen. Only the possibility ``$P$ is true'' remains. This is just case
analysis with two cases: $P$ true or $P$ false.
\end{proofbox}
\begin{tipbox}
Start with the ``tightest'' piece of information -- the one that allows the fewest possibilities. For instance, ``Jack and Kate do
\emph{not} sit next to each other'' at a $5$-seat table means there is \emph{exactly one} person between them on one side. Fewer
possibilities means fewer cases to check!
\end{tipbox}

% ---------------------------------------------------------------------
\needspace{9\baselineskip}\shead{Ordering; ``must be true'' and ``could be true''}
\begin{knowbox}
\textbf{1. Use symbols.} Translate each sentence into an inequality: ``Ryan runs faster than Henry'' becomes $R > H$;
``Kayla ate less than Max'' becomes $K < M$ (or $M > K$).

\textbf{2. Transitivity.} If $A > B$ and $B > C$, then $A > C$. This lets us chain inequalities together.

\textbf{3. Diagrams.} Draw each person as a point on a \emph{vertical axis}: whoever is ``more'' goes higher; or draw an arrow from
the ``lesser'' person to the ``greater'' one. Two people who are not joined by a chain of arrows in the same direction
\emph{cannot be compared} yet.

\textbf{4. Position in a line.} If someone is $k$-th from the front and $m$-th from the back of a line, the line has $k + m - 1$ people.

\textbf{5. ``Must be true'' and ``could be true''.}
\begin{itemize}[leftmargin=1.4em,itemsep=1pt,topsep=2pt]
  \item A statement that \textbf{must be true} (is always true): true in \emph{every} situation that satisfies the problem.
  \item A statement that \textbf{could be true}: true in \emph{at least one} situation that satisfies the problem.
  \item To refute a ``must be true'' claim, \emph{one counterexample} is enough: a situation that satisfies the problem but in which
  the claim is false.
\end{itemize}

\textbf{6. ``Every A is a B''.} From this we may conclude: ``whoever is \emph{not} a B is \emph{not} an A''. But we may \emph{not}
conclude ``whoever is a B is an A''.
\end{knowbox}
\begin{figC}
  \begin{scope}
  \foreach \y/\n in {2.4/Oscar,1.8/Tom,1.2/Ryan,0.6/Felix,0/Henry}{\fill[tmBlue] (0,\y) circle (2.2pt); \node[right,font=\small] at (0.12,\y){\n};}
  \draw[->,Navy,line width=0.8pt] (-0.35,-0.15)--(-0.35,2.6) node[above,font=\scriptsize,text=Slate]{faster};
  \end{scope}
  \begin{scope}[xshift=5.2cm]
  \draw[tmBlue,line width=1pt,fill=tmBlue!8] (0.9,1.1) ellipse (1.7 and 1.15);
  \draw[tmOrange,line width=1pt,fill=tmOrange!15] (0.55,1.0) ellipse (0.75 and 0.6);
  \node[font=\scriptsize,text=tmOrange!70!black] at (0.55,1.0){chess club};
  \node[font=\scriptsize,text=tmBlue!80!black] at (1.75,1.85){can swim};
  \fill[AnsRed] (3.2,1.1) circle (2pt) node[right,font=\scriptsize]{cannot swim};
  \end{scope}
\end{figC}
\begin{proofbox}
Why does being $k$-th from the front and $m$-th from the back give $k+m-1$ people? Counting from the front up to that person gives
$k$ people, counting from that person to the back gives $m$ people. Adding gives $k+m$, but that person has been counted \emph{twice},
so we subtract $1$.

Why does ``every A is a B'' give ``not a B means not an A''? Look at the right-hand picture: the ``chess club'' oval lies
\emph{entirely} inside the ``can swim'' oval. A point outside the big oval is certainly outside the small one. A point inside the big
oval, however, may be inside or outside the small one.
\end{proofbox}
\begin{warnbox}
When you see ``$A > B$'' and ``$A > C$'', do not jump to conclusions about $B$ and $C$: they \emph{cannot be compared} yet. A claim
like ``$B > C$'' then only \emph{could} be true; it is not \emph{certainly} true.
\end{warnbox}

% ---------------------------------------------------------------------
\needspace{9\baselineskip}\shead{Shift ciphers and remainders}
\begin{knowbox}
\textbf{1. The Caesar cipher.} Arrange the $26$ letters A, B, C, \ldots, Z of the English alphabet in a circle (after Z comes A
again). To encrypt with \emph{shift} $k$: replace each letter by the letter $k$ places after it, clockwise. To decrypt: shift
\emph{back} $k$ places.

\textbf{2. Numbering the letters.} Number them A $=0$, B $=1$, \ldots, Z $=25$. To encrypt letter number $x$ with shift $k$:
compute $x + k$; if the result is $\ge 26$, subtract $26$ (that is, take the \emph{remainder} of $x+k$ on division by $26$).
\begin{center}
\renewcommand{\arraystretch}{1.1}\setlength{\tabcolsep}{2.6pt}\small
\begin{tabular}{|*{13}{c|}}
\hline
\rowcolor{tmBlue!10}A&B&C&D&E&F&G&H&I&J&K&L&M\\ \hline
0&1&2&3&4&5&6&7&8&9&10&11&12\\ \hline
\rowcolor{tmBlue!10}N&O&P&Q&R&S&T&U&V&W&X&Y&Z\\ \hline
13&14&15&16&17&18&19&20&21&22&23&24&25\\ \hline
\end{tabular}
\end{center}
For example, with shift $4$: Y $= 24$, $24 + 4 = 28 \ge 26$, $28 - 26 = 2 =$ C. So Y $\to$ C.

\textbf{3. Shifting several times.} Shifting by $a$ and then by $b$ is the same as shifting once by $a + b$. Shifting by $26$ brings
every letter back to itself, so only the \emph{remainder} of the number of steps on division by $26$ matters.

\textbf{4. Days of the week ``go round'' too.} After Sunday comes Monday again: the days repeat in a cycle of $7$. Which day it will
be ``in $n$ days'' depends only on the remainder of $n$ on division by $7$.
\end{knowbox}
\begin{figC}
  \node at (0,0){\lgaWheel{4}};
  \node[align=left,font=\small,anchor=west] at (2.5,0){Shift $4$:\\ A $\to$ E, \ B $\to$ F, \ldots\\ W $\to$ A, \ X $\to$ B,\\ Y $\to$ C, \ Z $\to$ D.};
\end{figC}
\begin{proofbox}
Why does a shift of $26$ bring us back to the start? The circle has exactly $26$ letters, so $26$ steps is one full turn. Likewise,
a shift of $30 = 26 + 4$ is one full turn followed by $4$ more steps: exactly the same as a shift of $4$. And $4$ is precisely the
remainder when $30$ is divided by $26$.
\end{proofbox}
\begin{warnbox}
The most common mistake is forgetting to ``wrap around'' at the end of the alphabet: W shifted by $4$ does \emph{not} give some
strange symbol after Z, but A (W $\to$ X $\to$ Y $\to$ Z $\to$ A). The second mistake: decrypting by shifting \emph{forwards}
instead of \emph{backwards}.
\end{warnbox}

% ---------------------------------------------------------------------
\needspace{9\baselineskip}\shead{Magic squares}
\begin{knowbox}
\textbf{1. Definition.} A $3\times 3$ magic square is a grid of $9$ cells filled with $9$ different numbers so that the sums of the
numbers in every row, every column and both diagonals are equal. This common sum is called the \textbf{magic sum}, written $S$.

\textbf{2. The magic sum.} The three rows cover the whole grid, each cell exactly once, so
\[ 3S = \text{sum of all } 9 \text{ numbers}, \qquad S = \frac{\text{sum of all } 9 \text{ numbers}}{3}. \]
Example: for the numbers $1, 2, \ldots, 9$ the total is $45$, so $S = 15$.

\textbf{3. The centre number.} In a $3\times 3$ magic square: \ \textbf{centre number $= S \div 3$} (the average of the $9$ numbers).
For the numbers $1, \ldots, 9$ the centre is $5$.

\textbf{4. Quick filling trick.} Once you know $S$ and a row (column, diagonal) already has two numbers, the missing number is
$S -$ (sum of the two known numbers).
\end{knowbox}
\begin{figC}
  \node at (0,0){\lgaMagic{$2$}{$7$}{$6$}{$9$}{$5$}{$1$}{$4$}{$3$}{$8$}};
  \begin{scope}[xshift=4.2cm]
  \fill[tmOrange!25] (-1.35,-0.45) rectangle (1.35,0.45);
  \fill[tmOrange!25] (-0.45,-1.35) rectangle (0.45,1.35);
  \draw[tmPurple!70,line width=1.4pt] (-1.25,1.25)--(1.25,-1.25) (-1.25,-1.25)--(1.25,1.25);
  \draw[Navy,line width=0.9pt] (-1.35,-1.35) grid[step=0.9] (1.35,1.35);
  \node[font=\bfseries] at (0,0){$c$};
  \node[font=\footnotesize,text=Slate,align=center] at (0,-1.75){$4$ lines through the centre};
  \end{scope}
\end{figC}
\begin{proofbox}
Why is the centre number equal to $S \div 3$? Exactly $4$ lines pass through the centre: the middle row, the middle column and the
two diagonals (right-hand picture). Adding the sums of these $4$ lines gives $4S$. In that sum the centre $c$ is counted $4$ times,
while every other cell is counted exactly once (check on the picture: each outer cell lies on exactly one of the four lines). Hence
\[ 4S = (\text{sum of the } 9 \text{ numbers}) + 3c = 3S + 3c. \]
Subtract $3S$ from both sides: $S = 3c$, that is, $c = S \div 3$.
\end{proofbox}
\begin{tipbox}
Magic squares may use negative numbers too. With the numbers $-6, -5, \ldots, 2$: the total is $-18$, so $S = -6$ and the centre is
$-2$. The formula works just the same!
\end{tipbox}

% ---------------------------------------------------------------------
\needspace{9\baselineskip}\shead{A balance scale is an equation}
\begin{knowbox}
\textbf{1. Balanced scale $=$ equality.} If a scale balances, the total masses on the two pans are equal. Giving each kind of object
a letter for its mass, we can write an equality (an equation).

\textbf{2. Removing identical objects.} Removing identical objects from \emph{both} pans keeps the scale balanced (subtract the same
number from both sides).

\textbf{3. Substitution.} If a group of objects $A$ balances a group $B$, then on any other scale we may replace $A$ by $B$ and the
scale stays balanced.

\textbf{4. Reduce to one unknown.} Let the mass of \emph{one} kind of object be $x$, then use the equalities to express the mass of
every other object in terms of $x$. Comparing two groups of objects then becomes comparing two expressions in $x$.
\end{knowbox}
\begin{figC}
  \node at (0,0){\lgaScale{\lgaCir\,\lgaCir\,\lgaCir\,\lgaCir}{\lgaTri\,\lgaTri\,\lgaTri}};
  \node[font=\small,anchor=west,align=left] at (2.7,0){$4O = 3T$. \ Let $O = x$: \ $3T = 4x$, \ $T = \dfrac{4x}{3}$.};
\end{figC}
\begin{proofbox}
Why may we do this? A balanced scale means that the two sides of the equation $4O = 3T$ are equal. Dividing both sides by the same
non-zero number still gives a true equation: $T = \dfrac{4O}{3}$. So however many grams a circle really weighs, a triangle always
weighs $\dfrac{4}{3}$ of a circle. For instance, if $x = 30$ g then $T = 40$ g; check: $4\cdot 30 = 120 = 3\cdot 40$.
\end{proofbox}
\begin{warnbox}
Never remove an object from \emph{one} pan without removing an identical object from the other. And only remove objects that are
\emph{identical} (have the same mass).
\end{warnbox}

% ---------------------------------------------------------------------
\needspace{9\baselineskip}\shead{Clever counting: sums on figures, dice, fair games}
\begin{knowbox}
\textbf{1. How many times is each number counted?} When numbers are written in circles along the sides of a figure and every side
has the same sum $S$, add up the sums of \emph{all the sides}. A number at a \emph{corner} belongs to two sides, so it is counted
twice; a number in the \emph{middle of a side} is counted once. For a triangle:
\[ 3S = (\text{sum of all the numbers}) + (\text{sum of the three corner numbers}). \]

\textbf{2. Lower and upper bounds.} The sum of the three corners cannot be smaller than the sum of the three smallest numbers, nor
larger than the sum of the three largest. This pins down $S$; then check each possible value.

\textbf{3. Dice.} On a standard die, \emph{opposite faces add up to $7$}: $1$ -- $6$, $2$ -- $5$, $3$ -- $4$. The six faces add up to
$1+2+3+4+5+6 = 21$. Therefore: sum of the \emph{hidden} faces $= 21 -$ sum of the visible faces.

\textbf{4. Fair games.} Rolling a die has $6$ equally likely outcomes. A game ``you win if \ldots'' is \emph{fair} when the chance of
winning equals the chance of losing, that is, there are exactly $3$ winning and $3$ losing outcomes.
\end{knowbox}
\begin{figC}
  \node at (0,0){\lgaTriSix{$a$}{$b$}{$c$}{$x$}{$y$}{$z$}};
  \node[font=\small,anchor=west,align=left] at (1.8,0.1){left side: $a + x + b$\\ base: \hspace{6.4mm}$b + y + c$\\ right side: $c + z + a$\\[2pt] \textcolor{tmBlue!80!black}{$a, b, c$ appear twice.}};
  \node at (7.4,0){\lgaCube{3}{6}{2}};
  \node[font=\footnotesize,text=Slate,align=center] at (7.6,-1.0){visible $3 + 6 + 2 = 11$\\ hidden $21 - 11 = 10$};
\end{figC}
\begin{proofbox}
Why is $3S = (\text{sum of all the numbers}) + (a + b + c)$? Add the three sides:
$(a + x + b) + (b + y + c) + (c + z + a) = (a + b + c + x + y + z) + (a + b + c)$.
Each middle number $x, y, z$ appears once; each corner number $a, b, c$ appears twice, so $a + b + c$ must be added once more.
\end{proofbox}
\begin{warnbox}
For a die standing on a table, the \emph{bottom} face is hidden as well. Do not forget it! The safest method is to use ``sum of the
$6$ faces $= 21$'' and subtract the visible faces.
\end{warnbox}

% ---------------------------------------------------------------------
\needspace{9\baselineskip}\shead{Latin squares and sudoku}
\begin{knowbox}
\textbf{1. $4\times 4$ Latin square.} Write the digits $1, 2, 3, 4$ in a $4\times 4$ grid so that every row and every column contains
all four digits (no repeats). A \textbf{$4\times 4$ sudoku} adds one more condition: each $2\times 2$ box (thick border) also contains
all four digits.

\textbf{2. Three ``regions''.} Each cell belongs to exactly three regions: one row, one column, one $2\times 2$ box. The digit in the
cell must be \emph{different} from every digit already present in all three regions.

\textbf{3. Two ways of reasoning.}
\begin{itemize}[leftmargin=1.4em,itemsep=1pt,topsep=2pt]
  \item \emph{Elimination in one cell:} cross out every digit already in the cell's row, column and box; if exactly one digit is left,
  write it in.
  \item \emph{Finding a place for a digit:} if, in some region, digit $d$ can go in only one cell, then $d$ must go there.
\end{itemize}

\textbf{4. Write ``candidates''.} If a cell still has two possibilities, write both in small (e.g.\ $\{1, 4\}$) and move on;
information elsewhere will settle it.
\end{knowbox}
\begin{figC}
  \node at (0,0){\lgaFour{1/1/$1$,3/1/$3$,2/2/$2$,4/4/$4$}};
  \node[font=\small,anchor=west,align=left] at (1.8,0){Cell in row $1$, column $2$: row $1$ already has $1, 3$;\\ the $2\times 2$ box already has $1, 2$.\\ The only digit left is $4$.};
\end{figC}
\begin{warnbox}
Do not ``guess'' and write the digit straight in. Every digit you write must have a reason (all other possibilities are ruled out).
If you really must try something, write ``suppose'' clearly and be ready to go back when you meet a contradiction.
\end{warnbox}
\begin{parentbox}
Logic puzzles are great to do together with your child: read the problem aloud, let your child retell it in their own words, then
draw a diagram or a table together. When your child gives an answer, do not say at once whether it is right or wrong, but ask:
``Have you checked it against \emph{each} condition of the problem?''. The habit of checking is worth more than any trick.
\end{parentbox}

% =====================================================================
% WORKED EXAMPLES
% =====================================================================
\begin{exbox}{Example 1 \ (Truth-tellers and liars)}
On an island, every person either always tells the truth or always lies. You meet two people, Peter and Quinn.
Peter says: ``Quinn is a liar.'' Quinn says: ``Peter and I are the same type of person.''
Who tells the truth and who lies?
\guide
\step{1}{Case: Peter tells the truth}
Then Peter's statement is true: Quinn is a liar. So Quinn's statement must be false. Quinn says ``we are the same type''; in fact
Peter is truthful and Quinn lies -- \emph{different} types, so Quinn's statement is indeed false. \textbf{No contradiction.}
\step{2}{Case: Peter lies}
Then Peter's statement is false: Quinn is \emph{not} a liar, so Quinn tells the truth. Then Quinn's statement must be true: ``Peter and
Quinn are the same type''. But Peter lies and Quinn tells the truth -- different types. So Quinn's statement is false after all.
\textbf{Contradiction!} This case is ruled out.
\step{3}{Summary in a table}
\begin{center}
\begin{tabular}{|c|c|c|c|c|}
\hline
\rowcolor{tmBlue!10} Peter & Quinn & Peter's statement & Quinn's statement & Conclusion \\ \hline
truthful & liar & true (fine) & false (fine) & \textcolor{tmGreen!70!black}{\textbf{OK}} \\ \hline
liar & truthful & false (fine) & false (impossible!) & \textcolor{AnsRed}{ruled out} \\ \hline
\end{tabular}
\end{center}
(The cases ``both truthful'' and ``both liars'' are ruled out straight away by Peter's statement: if Peter is truthful then Quinn lies;
if Peter lies then Quinn is truthful.)
\ans{Peter tells the truth, Quinn lies.}
\end{exbox}

\begin{exbox}{Example 2 \ (Each person is half right)}
Three teams, Red, Blue and Yellow, play a tug-of-war; no two teams share a place. Before the match, three fans predict:
\begin{itemize}[leftmargin=1.4em,itemsep=0pt,topsep=2pt]
\item Grace: ``Red first, Blue second.'' \quad Jack: ``Blue first, Yellow second.'' \quad Lucy: ``Blue third, Yellow first.''
\end{itemize}
In the end, each person got exactly one part right. Find the places of the three teams.
\guide
\step{1}{Pick one part to split into cases}
Each prediction has exactly one true part. Look at Grace's: either ``Red first'' is true or ``Blue second'' is true.
\step{2}{Case 1: ``Blue second'' is true (and ``Red first'' is false)}
If Blue is second, Blue is not first, so Jack's part ``Blue first'' is false; hence his other part ``Yellow second'' must be true. But
Blue is already second, so Yellow cannot be second as well. \textbf{Contradiction}, ruled out.
\step{3}{Case 2: ``Red first'' is true (and ``Blue second'' is false)}
Blue is not first (Red is) and not second, so \textbf{Blue is third}. That leaves Yellow second.
\step{4}{Check all three people}
\begin{center}
\begin{tabular}{|l|c|c|}
\hline
\rowcolor{tmBlue!10} Person & First part & Second part \\ \hline
Grace & Red first: \textcolor{tmGreen!70!black}{true} & Blue second: \textcolor{AnsRed}{false} \\ \hline
Jack & Blue first: \textcolor{AnsRed}{false} & Yellow second: \textcolor{tmGreen!70!black}{true} \\ \hline
Lucy & Blue third: \textcolor{tmGreen!70!black}{true} & Yellow first: \textcolor{AnsRed}{false} \\ \hline
\end{tabular}
\end{center}
Each person has exactly one part right.
\ans{Red first, Yellow second, Blue third.}
\end{exbox}

\begin{exbox}{Example 3 \ (Ordering -- must or could?)}
There are four coloured pencils, red, blue, yellow and purple, all of different lengths. We know: the red pencil is longer than the blue
one; the yellow pencil is shorter than the blue one; the purple pencil is longer than the yellow one but shorter than the red one.\\
a) Which pencil is certainly the longest? Which is certainly the shortest?\\
b) Is the statement ``The purple pencil is the second longest'' certainly true, possibly true, or certainly false?
\guide
\step{1}{Use symbols}
Call the lengths $R, B, Y, P$. We are told: $R > B$; \ $B > Y$; \ $P > Y$; \ $R > P$.
\step{2}{Draw an arrow diagram (arrows point from the shorter pencil to the longer one)}
\begin{figC}
  \node[draw=tmOrange,rounded corners=3pt,line width=0.8pt] (V) at (0,0){Yellow};
  \node[draw=tmOrange,rounded corners=3pt,line width=0.8pt] (X) at (2.2,0.7){Blue};
  \node[draw=tmOrange,rounded corners=3pt,line width=0.8pt] (T) at (2.2,-0.7){Purple};
  \node[draw=tmOrange,rounded corners=3pt,line width=0.8pt] (D) at (4.4,0){Red};
  \draw[->,Navy,line width=0.9pt] (V)--(X); \draw[->,Navy,line width=0.9pt] (V)--(T);
  \draw[->,Navy,line width=0.9pt] (X)--(D); \draw[->,Navy,line width=0.9pt] (T)--(D);
  \draw[<->,dashed,AnsRed!70] (X)--(T) node[midway,right,font=\scriptsize]{?};
\end{figC}
\step{3}{Answer part a}
From $R > B > Y$ and $R > P > Y$: the red pencil is longer than all three others and the yellow pencil is shorter than all three others.
So \textbf{red is certainly the longest} and \textbf{yellow is certainly the shortest}.
\step{4}{Answer part b: find two situations}
No condition compares Blue with Purple. Try two sets of lengths (cm), both satisfying every condition:
\begin{itemize}[leftmargin=1.4em,itemsep=0pt,topsep=1pt]
\item Red $15$, Purple $13$, Blue $12$, Yellow $8$: purple is second longest -- the statement is \emph{true};
\item Red $15$, Blue $12$, Purple $10$, Yellow $8$: purple is third longest -- the statement is \emph{false}.
\end{itemize}
So ``The purple pencil is the second longest'' is only \textbf{possibly true}, not certainly true.
\ans{a) Red is the longest, yellow the shortest (certainly). \ b) Possibly true (not certain).}
\end{exbox}

\begin{exbox}{Example 4 \ (Shift cipher -- finding the shift)}
Tom encrypts a word with a Caesar cipher and gets \textbf{JVKL}. He tells us that with this encryption the letter A becomes H.\\
a) Find the original word. \ b) After this shift, what is the smallest number of further steps (in the same direction) that brings every
letter back to its original letter?
\guide
\step{1}{Find the shift}
A $= 0$, H $= 7$. A becomes H, so the shift is $k = 7$.
\step{2}{Decrypt: shift back $7$ (subtract $7$; if the result is negative, add $26$)}
\begin{center}
\begin{tabular}{|c|c|c|c|c|}
\hline
\rowcolor{tmBlue!10} Cipher letter & J & V & K & L \\ \hline
Number & $9$ & $21$ & $10$ & $11$ \\ \hline
Subtract $7$ & $2$ & $14$ & $3$ & $4$ \\ \hline
Original letter & \textbf{C} & \textbf{O} & \textbf{D} & \textbf{E} \\ \hline
\end{tabular}
\end{center}
The original word is \textbf{CODE}.
\begin{warnbox}Subtracting $7$ does not always give a negative number; we add $26$ (wrap around past A) only when the number of the
cipher letter is less than $7$. Here all four letters have numbers $\ge 7$, so no wrapping is needed.\end{warnbox}
\step{3}{Part b}
A total shift of $26$ steps goes once round the circle. We have shifted $7$ steps, so we need $26 - 7 = 19$ more steps.
\verify{Encrypt CODE again with shift $7$: C$(2)\to 9 =$ J, O$(14)\to 21 =$ V, D$(3)\to 10 =$ K, E$(4)\to 11 =$ L. We get JVKL.}
\ans{a) CODE. \ b) $19$ steps.}
\end{exbox}

\begin{exbox}{Example 5 \ (Completing a magic square)}
Place the even numbers $2, 4, 6, \ldots, 18$ (each once) in the grid so that every row, column and diagonal has the same sum.
\begin{center}\lgaMagic{$4$}{}{$12$}{}{}{}{}{$6$}{}\end{center}
\guide
\step{1}{Find the magic sum $S$}
Sum of the nine numbers: $2 + 4 + \cdots + 18 = 2\times(1 + 2 + \cdots + 9) = 2\times 45 = 90$. So $S = 90 \div 3 = 30$.
\step{2}{The centre number}
\begin{recallbox}In a $3\times 3$ magic square the centre number equals $S \div 3$.\end{recallbox}
Centre $= 30 \div 3 = 10$.
\step{3}{Fill in the lines that already have two numbers}
\begin{itemize}[leftmargin=1.4em,itemsep=0pt,topsep=1pt]
\item Top row: middle cell $= 30 - 4 - 12 = 14$.
\item Diagonal $4$ -- $10$ -- ?: bottom-right corner $= 30 - 4 - 10 = 16$.
\item Diagonal $12$ -- $10$ -- ?: bottom-left corner $= 30 - 12 - 10 = 8$.
\item Left column: $30 - 4 - 8 = 18$; \ right column: $30 - 12 - 16 = 2$.
\end{itemize}
\begin{center}\lgaMagic{$4$}{\textcolor{tmGreen!60!black}{$14$}}{$12$}{\textcolor{tmGreen!60!black}{$18$}}{\textcolor{AnsRed}{$10$}}{\textcolor{tmGreen!60!black}{$2$}}{\textcolor{tmGreen!60!black}{$8$}}{$6$}{\textcolor{tmGreen!60!black}{$16$}}\end{center}
\verify{Bottom row: $8 + 6 + 16 = 30$; middle row: $18 + 10 + 2 = 30$; middle column: $14 + 10 + 6 = 30$. Every even number from $2$ to $18$ is used exactly once.}
\ans{Top row $4, 14, 12$; middle row $18, 10, 2$; bottom row $8, 6, 16$.}
\end{exbox}

\begin{exbox}{Example 6 \ (A balance-scale problem)}
Both scales below are balanced; identical shapes have equal masses. How many circles does one square weigh?
\begin{center}
\lgaScale{\lgaSq\,\lgaSq}{\lgaTri\,\lgaTri\,\lgaTri}\hspace{1cm}\lgaScale{\lgaTri}{\lgaCir\,\lgaCir}
\end{center}
\guide
\step{1}{Write equalities}
Let the masses of a square, a triangle and a circle be $V, T, O$. Scale 1: $2V = 3T$. \ Scale 2: $T = 2O$.
\step{2}{Substitute}
On scale 1, replace each triangle by two circles: $2V = 3\times 2O = 6O$.
\step{3}{Halve both pans}
$2V = 6O$, so $V = 3O$.
\verify{Take $O = 1$, $T = 2$, $V = 3$: scale 1: $2\times 3 = 6 = 3\times 2$; scale 2: $2 = 2\times 1$. Correct.}
\ans{One square weighs as much as $3$ circles.}
\end{exbox}

\begin{exbox}{Example 7 \ (Sums along the sides of a triangle)}
Place the numbers $1, 2, 3, 4, 5, 6$ in the six circles (three at the corners, three in the middles of the sides) so that the three
numbers on each side add up to the same number $S$. Find all possible values of $S$.
\guide
\step{1}{Count how many times each number is counted}
\begin{recallbox}When adding the three side sums, corner numbers are counted twice and middle numbers once.\end{recallbox}
$3S = (1 + 2 + \cdots + 6) + (\text{sum of the three corners}) = 21 + (\text{sum of the three corners})$.
\step{2}{Pin it down}
The smallest possible corner sum is $1 + 2 + 3 = 6$, the largest is $4 + 5 + 6 = 15$. So $27 \le 3S \le 36$, that is, $9 \le S \le 12$.
\step{3}{Check each value with an actual arrangement}
\begin{figC}
  \node at (0,0){\lgaTriSix{$1$}{$2$}{$3$}{$6$}{$4$}{$5$}}; \node[font=\small] at (0,-1.6){$S = 9$};
  \node at (3.3,0){\lgaTriSix{$1$}{$3$}{$5$}{$6$}{$2$}{$4$}}; \node[font=\small] at (3.3,-1.6){$S = 10$};
  \node at (6.6,0){\lgaTriSix{$2$}{$4$}{$6$}{$5$}{$1$}{$3$}}; \node[font=\small] at (6.6,-1.6){$S = 11$};
  \node at (9.9,0){\lgaTriSix{$4$}{$5$}{$6$}{$3$}{$1$}{$2$}}; \node[font=\small] at (9.9,-1.6){$S = 12$};
\end{figC}
For instance, with $S = 10$: $1 + 6 + 3 = 10$, \ $3 + 2 + 5 = 10$, \ $5 + 4 + 1 = 10$.
\ans{$S$ can be $9$, $10$, $11$ or $12$.}
\end{exbox}

\begin{exbox}{Example 8 \ ($4\times 4$ sudoku)}
Complete the sudoku: every row, every column and every thick-bordered $2\times 2$ box must contain the digits $1, 2, 3, 4$.
\begin{center}\lgaFour{1/1/$2$,2/1/$1$,3/2/$1$,2/3/$4$,4/4/$3$}\end{center}
\guide
\step{1}{Top-left box}
It already has $2, 1$ in row $1$, so row $2$, columns $1$ and $2$, are $\{3, 4\}$. Column $2$ already has a $4$ (in row $3$), so
row $2$ column $2$ is $3$ and row $2$ column $1$ is $4$.
\step{2}{Row $2$ and the top-right box}
Row $2$: $4, 3, 1, \_$ $\Rightarrow$ the last cell is $2$. Row $1$, columns $3$ and $4$, are $\{3, 4\}$; column $4$ already has a $3$
(row $4$), so row $1$ column $4$ is $4$ and row $1$ column $3$ is $3$.
\step{3}{Column $2$ and the bottom-left box}
Column $2$: $1, 3, 4, \_$ $\Rightarrow$ row $4$ column $2$ is $2$. Column $1$ is still missing $\{1, 3\}$; row $4$ already has a $3$, so
row $4$ column $1$ is $1$ and row $3$ column $1$ is $3$.
\step{4}{The last two rows}
Row $3$: $3, 4, \_, \_$ needs $\{1, 2\}$; column $3$ already has a $1$ (row $2$), so row $3$ column $3$ is $2$ and column $4$ is $1$. Row $4$: $1, 2, \_, 3$ $\Rightarrow$ the empty cell is $4$.
\begin{center}\lgaFour{1/1/$2$,2/1/$1$,3/1/\textcolor{tmGreen!60!black}{$3$},4/1/\textcolor{tmGreen!60!black}{$4$},1/2/\textcolor{tmGreen!60!black}{$4$},2/2/\textcolor{tmGreen!60!black}{$3$},3/2/$1$,4/2/\textcolor{tmGreen!60!black}{$2$},1/3/\textcolor{tmGreen!60!black}{$3$},2/3/$4$,3/3/\textcolor{tmGreen!60!black}{$2$},4/3/\textcolor{tmGreen!60!black}{$1$},1/4/\textcolor{tmGreen!60!black}{$1$},2/4/\textcolor{tmGreen!60!black}{$2$},3/4/\textcolor{tmGreen!60!black}{$4$},4/4/$3$}\end{center}
\verify{Every row, column and $2\times 2$ box contains all of $1, 2, 3, 4$.}
\ans{Row $1$: $2\,1\,3\,4$; row $2$: $4\,3\,1\,2$; row $3$: $3\,4\,2\,1$; row $4$: $1\,2\,4\,3$.}
\end{exbox}

\clearpage
\dividerbar{Core practice}

% ---------------------------------------------------------------- Problem 1
\begin{probbox}{1}{\relax}
Mum has just bought a new pot of flowers. Three children guess what flower it is. Anna guesses: ``a red rose''. Ben guesses: ``a purple
daisy''. Chloe guesses: ``a red dahlia''. Mum says that each child got \emph{exactly one} of the two things right -- either the colour
or the kind of flower. The flower Mum bought is a
\choicesII{purple dahlia}{purple rose}{red dahlia}{yellow rose}{red daisy}
\end{probbox}
\probsol{
\textbf{Analysing the problem.} Each guess has two parts: a \emph{colour} and a \emph{kind of flower}. Each child has exactly one part
right: if one part is right, the other is wrong. We must find the real colour and kind.
\begin{recallbox}``Half-true'' statement: if we know one part is true, the other part is false, and vice versa.\end{recallbox}
\step{1}{The colour: is the flower red?}
\emph{Suppose the flower is not red.} Then Anna's part ``red'' is wrong, so her other part ``rose'' must be right. Likewise Chloe's part
``red'' is wrong, so ``dahlia'' must be right. The flower would be both a rose and a dahlia -- \textbf{impossible}. So
\textbf{the flower is red}.
\step{2}{Deduce the kind of flower}
The flower is red, so Anna has the colour right $\Rightarrow$ Anna's ``rose'' is wrong: it is not a rose.\\
Chloe also has the colour right $\Rightarrow$ Chloe's ``dahlia'' is wrong: it is not a dahlia.\\
Ben's ``purple'' is wrong (the flower is red), so Ben's ``daisy'' must be right: \textbf{a daisy}.
\step{3}{Check with a table}
\begin{center}
\begin{tabular}{|l|c|c|c|}
\hline
\rowcolor{tmBlue!10} Child & Colour guessed & Kind guessed & Parts right (real flower: red daisy) \\ \hline
Anna & red \textcolor{tmGreen!70!black}{(right)} & rose \textcolor{AnsRed}{(wrong)} & $1$ \\ \hline
Ben & purple \textcolor{AnsRed}{(wrong)} & daisy \textcolor{tmGreen!70!black}{(right)} & $1$ \\ \hline
Chloe & red \textcolor{tmGreen!70!black}{(right)} & dahlia \textcolor{AnsRed}{(wrong)} & $1$ \\ \hline
\end{tabular}
\end{center}
\begin{warnbox}Try option (B) ``purple rose'': Anna has the kind right (rose), Ben has the colour right (purple), but Chloe has
\emph{nothing} right -- against the conditions. You must check \emph{all three} children, not just one or two.\end{warnbox}
\verify{With ``red daisy'' each child has exactly one part right. Every other option gives some child $0$ or $2$ right parts: (A) purple dahlia -- Anna has $0$; (C) red dahlia -- Chloe has $2$; (D) yellow rose -- Ben has $0$.}
So the answer is \textbf{(E)}.
\ans{(E) red daisy}
}

% ---------------------------------------------------------------- Problem 2
\begin{probbox}{2}{\relax}
The $26$ letters of the English alphabet are written in order A, B, C, \ldots, Z clockwise around a circle. To \emph{encrypt} a message,
each letter is replaced by the letter $4$ places after it, clockwise (this is called a \emph{Caesar cipher}). For example, the word ZAP
is encrypted as DET. The word WIN is encrypted as
\choices{ALN}{ZLN}{AMR}{AMQ}{ZMQ}
\end{probbox}
\probsol{
\textbf{Analysing the problem.} The shift is $4$, forwards (clockwise). We must encrypt the letters W, I, N one by one.
\begin{recallbox}Number A $= 0$, \ldots, Z $= 25$. Encrypt: add $4$; if the result is $\ge 26$, subtract $26$ (going past Z brings you back to A).\end{recallbox}
\step{1}{Check our understanding with ZAP $\to$ DET}
Z$(25) + 4 = 29$, $29 - 26 = 3 =$ D; \ A$(0) + 4 = 4 =$ E; \ P$(15) + 4 = 19 =$ T. That is indeed DET, and we see Z ``wrap around'' to D.
\step{2}{Encrypt each letter of WIN}
\begin{center}
\begin{tabular}{|c|c|c|c|}
\hline
\rowcolor{tmBlue!10} Original letter & W & I & N \\ \hline
Number & $22$ & $8$ & $13$ \\ \hline
Add $4$ & $26 \to 26 - 26 = 0$ & $12$ & $17$ \\ \hline
Cipher letter & \textbf{A} & \textbf{M} & \textbf{R} \\ \hline
\end{tabular}
\end{center}
\begin{figC}
  \node at (0,0){\lgaWheel{4}};
  \node[align=left,font=\small,anchor=west] at (2.5,0){W $\to$ X $\to$ Y $\to$ Z $\to$ \textbf{A}\\ I $\to$ J $\to$ K $\to$ L $\to$ \textbf{M}\\ N $\to$ O $\to$ P $\to$ Q $\to$ \textbf{R}};
\end{figC}
\begin{warnbox}Each wrong option has at least one letter shifted by the wrong number of steps: (B) ZLN and (E) ZMQ shift W by only $3$
(stopping at Z and forgetting that after Z we carry on to A); (D) AMQ and (E) ZMQ shift N by only $3$; (A) ALN and (B) ZLN shift I by
$3$ and leave N unchanged. Count carefully: ``$4$ places after'' means taking exactly $4$ steps, not counting the starting letter.\end{warnbox}
\verify{Decrypt AMR by going back $4$ steps: A $\to$ Z $\to$ Y $\to$ X $\to$ W; \ M $\to$ I; \ R $\to$ N. We get WIN back.}
So the answer is \textbf{(C)}.
\ans{(C) AMR}
}

% ---------------------------------------------------------------- Problem 3
\begin{probbox}{3}{\relax}
Five friends, Ryan, Henry, Felix, Tom and Oscar, run a race; no two finish at the same time. Ryan runs faster than Henry and Felix. Henry runs
slower than Felix. Tom runs faster than Ryan but slower than Oscar. Who finishes fourth?
\choices{Felix}{Henry}{Oscar}{Ryan}{Tom}
\end{probbox}
\probsol{
\textbf{Analysing the problem.} Five runners and several speed comparisons. We must find the finishing order and the fourth runner.
Faster runners finish earlier.
\step{1}{Use symbols ($>$ means ``faster than'')}
\begin{itemize}[leftmargin=1.4em,itemsep=0pt,topsep=1pt]
\item ``Ryan is faster than Henry and Felix'': \ $R > H$ \ and \ $R > F$.
\item ``Henry is slower than Felix'': \ $F > H$.
\item ``Tom is faster than Ryan but slower than Oscar'': \ $O > T > R$.
\end{itemize}
\step{2}{Chain them together (transitivity)}
From $O > T > R$, \ $R > F$ and $F > H$ we get: \quad $O > T > R > F > H$.
\begin{figC}
  \foreach \x/\n/\k/\o in {0/Oscar/1/1st,2/Tom/2/2nd,4/Ryan/3/3rd,6/Felix/4/4th,8/Henry/5/5th}{
    \node[draw=tmBlue,fill=tmBlue!8,rounded corners=3pt,minimum width=1.3cm,line width=0.8pt] (n\k) at (\x,0){\n};
    \node[font=\scriptsize,text=Slate] at (\x,-0.5){\o};}
  \foreach \a/\b in {1/2,2/3,3/4,4/5}{\draw[->,Navy,line width=0.8pt] (n\a)--(n\b);}
  \node[font=\scriptsize,text=Slate] at (4,0.55){(arrow: ``faster than'')};
  \draw[AnsRed,line width=1.2pt,rounded corners=4pt] (5.25,-0.8) rectangle (6.75,0.35);
\end{figC}
\step{3}{Read off the result}
Finishing order: Oscar (1st), Tom (2nd), Ryan (3rd), \textbf{Felix (4th)}, Henry (5th).
\begin{warnbox}``Ryan is faster than Henry and Felix'' does not tell us which of Henry and Felix is faster -- we need ``Henry is slower than
Felix''. Do not skip any sentence!\end{warnbox}
\verify{Ryan (3rd) is faster than Henry (5th) and Felix (4th) \checkmark; Henry is slower than Felix \checkmark; Tom (2nd) is faster than Ryan and slower than Oscar \checkmark.}
So the answer is \textbf{(A)}.
\ans{(A) Felix}
}

% ---------------------------------------------------------------- Problem 4
\begin{probbox}{4}{\relax}
\begin{minipage}[t]{0.62\linewidth}
In a magic square, the numbers in every row, every column and both diagonals have the same sum. The magic square shown uses each of the
integers from $-6$ to $2$ exactly once. The value of $Y$ is
\choicesII{$-1$}{$0$}{$-6$}{$2$}{$-2$}
\end{minipage}\hfill
\begin{minipage}[t]{0.33\linewidth}\vspace{-2pt}\centering
\lgaMagic{$1$}{}{$Y$}{$-4$}{}{}{$-3$}{}{$-5$}
\end{minipage}
\end{probbox}
\probsol{
\textbf{Analysing the problem.} The nine numbers are $-6, -5, -4, -3, -2, -1, 0, 1, 2$. Four cells are known; we need the cell $Y$ in the
top-right corner.
\step{1}{Find the magic sum $S$}
Sum of the nine numbers: $(-6) + (-5) + (-4) + (-3) + (-2) + (-1) + 0 + 1 + 2$. Pair $(-2) + 2 = 0$, $(-1) + 1 = 0$; what is left is
$(-6) + (-5) + (-4) + (-3) = -18$. So the total is $-18$, and
\[ S = (-18) \div 3 = -6. \]
\verify{The left column is already complete: $1 + (-4) + (-3) = -6$. This agrees with $S = -6$.}
\step{2}{Find the centre number}
\begin{recallbox}In a $3\times 3$ magic square the centre number equals $S \div 3$.\end{recallbox}
Centre $= (-6) \div 3 = -2$. (Another way: the main diagonal gives $1 + c + (-5) = -6$, so $c = -6 - 1 + 5 = -2$.)
\step{3}{Find $Y$ using the other diagonal}
The diagonal through $Y$, the centre and the cell $-3$: \ $Y + (-2) + (-3) = -6$, so $Y = -6 + 2 + 3 = -1$.
\step{4}{Complete the grid to check}
Top row: middle cell $= -6 - 1 - (-1) = -6$. Middle row: right cell $= -6 - (-4) - (-2) = 0$. Bottom row: middle cell $= -6 - (-3) - (-5) = 2$.
\begin{center}\lgaMagic{$1$}{\textcolor{tmGreen!60!black}{$-6$}}{\textcolor{AnsRed}{$-1$}}{$-4$}{\textcolor{tmGreen!60!black}{$-2$}}{\textcolor{tmGreen!60!black}{$0$}}{$-3$}{\textcolor{tmGreen!60!black}{$2$}}{$-5$}\end{center}
\verify{Middle column: $-6 - 2 + 2 = -6$; right column: $-1 + 0 - 5 = -6$; middle row: $-4 - 2 + 0 = -6$. Each number from $-6$ to $2$ is used exactly once.}
\begin{warnbox}Be careful with minus signs: $-6 - (-4) = -6 + 4 = -2$. Subtracting a negative number means adding its opposite.\end{warnbox}
So the answer is \textbf{(A)}.
\ans{(A) $-1$}
}

% ---------------------------------------------------------------- Problem 5
\begin{probbox}{5}{\relax}
A game is called \emph{fair} if your chance of winning equals your chance of losing. You roll a fair six-sided die. How many of the
following four games are fair?
\begin{itemize}[leftmargin=1.6em,itemsep=0pt,topsep=2pt]
\item You win if you roll a $2$.
\item You win if you roll an even number.
\item You win if you roll a number less than $4$.
\item You win if you roll a number divisible by $3$.
\end{itemize}
\choices{$0$}{$1$}{$2$}{$3$}{$4$}
\end{probbox}
\probsol{
\textbf{Analysing the problem.} There are $6$ equally likely outcomes: $1, 2, 3, 4, 5, 6$. A game is fair when the number of winning
outcomes equals the number of losing outcomes, i.e.\ $3$ winning and $3$ losing outcomes.
\begin{recallbox}With a fair die all six faces are equally likely. Chance of winning $=$ (number of winning outcomes) $\div$ $6$.\end{recallbox}
\step{1}{List the winning outcomes of each game}
\begin{center}
\begin{tabular}{|l|c|c|c|c|}
\hline
\rowcolor{tmBlue!10} Game & Winning outcomes & Wins & Losses & Fair? \\ \hline
Roll a $2$ & $2$ & $1$ & $5$ & no \\ \hline
Even number & $2, 4, 6$ & $3$ & $3$ & \textbf{yes} \\ \hline
Less than $4$ & $1, 2, 3$ & $3$ & $3$ & \textbf{yes} \\ \hline
Divisible by $3$ & $3, 6$ & $2$ & $4$ & no \\ \hline
\end{tabular}
\end{center}
\step{2}{Count}
There are $2$ fair games: ``even number'' and ``less than $4$''.
\begin{warnbox}``Less than $4$'' does \emph{not} include $4$: only $1, 2, 3$. Counting $4$ by mistake gives $4$ winning outcomes -- not fair.
And ``divisible by $3$'' includes both $3$ and $6$ (do not forget $6$).\end{warnbox}
\verify{Game ``even number'': chance of winning $\dfrac{3}{6} = \dfrac12$, chance of losing also $\dfrac12$. The same holds for ``less than $4$''.}
So the answer is \textbf{(C)}.
\ans{(C) $2$}
}

% ---------------------------------------------------------------- Problem 6
\begin{probbox}{6}{\relax}
Five friends, Jack, Kate, Leo, Mia and Noah, sit on the $5$ chairs around a round table. Leo and Mia sit next to each other. Jack and Kate
do \emph{not} sit next to each other. The two people sitting on either side of Noah are
\choicesII{Jack and Leo}{Jack and Kate}{Kate and Mia}{Leo and Mia}{Mia and Jack}
\end{probbox}
\probsol{
\textbf{Analysing the problem.} Five people around a round table with $5$ chairs. Two conditions: Leo -- Mia are neighbours; Jack -- Kate
are not. We must find Noah's two neighbours.
\begin{recallbox}Around a round table we may fix one person's seat (rotating the table does not change who sits next to whom). With $5$
chairs, each person is not next to exactly $2$ people -- the two who are ``one seat away''.\end{recallbox}
\step{1}{Start with the tightest condition: Jack and Kate are not neighbours}
Number the chairs $1$ to $5$ around the table. Put Jack on chair $1$. The chairs next to chair $1$ are chairs $2$ and $5$, so Kate sits on
chair $3$ or chair $4$. These two cases are mirror images of each other, so we consider Kate on chair $3$.
\step{2}{Seat Leo and Mia}
Chairs $2$, $4$, $5$ are free. Leo and Mia need two chairs \emph{next to each other}. Of these three chairs, only chairs $4$ and $5$ are
neighbours (chair $2$ is next to chairs $1$ and $3$, both taken). So Leo and Mia sit on chairs $4$ and $5$.
\step{3}{Noah takes the remaining chair}
Noah sits on chair $2$, between chair $1$ (Jack) and chair $3$ (Kate).
\begin{figC}
  \foreach \i in {1,...,5}{\coordinate (S\i) at ({90-(\i-1)*72}:1.35);}
  \draw[Navy!40,line width=5pt] (0,0) circle (0.75);
  \foreach \i/\n/\c in {1/Jack/tmBlue,2/Noah/AnsRed,3/Kate/tmBlue,4/Leo/tmGreen,5/Mia/tmGreen}{
    \node[circle,draw=\c,fill=\c!10,line width=0.9pt,minimum size=0.95cm,font=\scriptsize] at (S\i){\n};
    \node[font=\scriptsize,text=Slate] at ({90-(\i-1)*72}:2.4){chair \i};}
\end{figC}
(If Kate is on chair $4$, it is similar: Leo and Mia are on chairs $2$, $3$ and Noah is on chair $5$ -- still between Jack and Kate.)
\begin{proofbox}Why is there no other arrangement? Jack and Kate are not neighbours, so between them there is exactly one chair on one
side and two chairs on the other side. Leo -- Mia need two adjacent chairs, so they must take the side with two chairs. The single chair
between Jack and Kate belongs to Noah.\end{proofbox}
\verify{Leo -- Mia are neighbours \checkmark; Jack (chair $1$) and Kate (chair $3$) are one chair apart \checkmark.}
So the answer is \textbf{(B)}.
\ans{(B) Jack and Kate}
}

% ---------------------------------------------------------------- Problem 7
\begin{probbox}{7}{\relax}
\begin{minipage}[t]{0.62\linewidth}
On each die, the numbers of pips on opposite faces always add up to $7$. Two dice stand on a table as shown. The total number of pips on
all the faces of the two dice that \emph{cannot be seen} is
\choicesII{$14$}{$20$}{$21$}{$24$}{$30$}
\end{minipage}\hfill
\begin{minipage}[t]{0.34\linewidth}\vspace{2pt}\centering
\lgaCube{3}{6}{2}\hspace{0.3cm}\lgaCube{1}{5}{4}
\end{minipage}
\end{probbox}
\probsol{
\textbf{Analysing the problem.} Each die has $6$ faces; we see $3$ of them (top, front, right) and $3$ are hidden (bottom, back, left).
We need the total number of pips on the $6$ hidden faces of the two dice.
\begin{recallbox}The six faces of a die carry $1 + 2 + 3 + 4 + 5 + 6 = 21$ pips in total.\end{recallbox}
\step{1}{Read the visible faces}
First die: top $3$, front $6$, right $2$ $\Rightarrow$ total $3 + 6 + 2 = 11$.\\
Second die: top $1$, front $5$, right $4$ $\Rightarrow$ total $1 + 5 + 4 = 10$.
\step{2}{Compute the hidden faces}
First die: $21 - 11 = 10$. \quad Second die: $21 - 10 = 11$. \quad Total: $10 + 11 = 21$.
\step{3}{Second method: use opposite faces}
Each hidden face is opposite a visible face, so it has $7 - (\text{visible face})$ pips.
First die: $(7 - 3) + (7 - 6) + (7 - 2) = 4 + 1 + 5 = 10$. Second die: $(7 - 1) + (7 - 5) + (7 - 4) = 6 + 2 + 3 = 11$.
\begin{proofbox}Why is each hidden face opposite exactly one visible face? The three visible faces (top, front, right) are pairwise
adjacent, so no two of them are opposite. The three faces opposite them (bottom, back, left) are exactly the three hidden faces.\end{proofbox}
\begin{warnbox}Some students count only the two hidden faces ``at the back'' and forget the bottom face resting on the table. The bottom
face cannot be seen either!\end{warnbox}
\verify{All faces of both dice total $21 + 21 = 42$; visible faces $11 + 10 = 21$; hidden faces $42 - 21 = 21$.}
So the answer is \textbf{(C)}.
\ans{(C) $21$}
}

% ---------------------------------------------------------------- Problem 8
\begin{probbox}{8}{\relax}
Chris lies on Fridays, Saturdays and Sundays and tells the truth on the other days. Mark lies on Tuesdays, Wednesdays and Thursdays and
tells the truth on the other days. On which day of the week could both of them say: ``Tomorrow I will lie''?
\choices{Monday}{Thursday}{Friday}{Sunday}{Tuesday}
\end{probbox}
\probsol{
\textbf{Analysing the problem.} A person can say ``Tomorrow I will lie'' in only two situations:
\begin{itemize}[leftmargin=1.4em,itemsep=0pt,topsep=1pt]
\item they tell the \emph{truth} today and the statement is \emph{true}: truthful today, lying tomorrow;
\item they \emph{lie} today and the statement is \emph{false}: lying today, truthful tomorrow.
\end{itemize}
In short: the person can say it when \textbf{today and tomorrow are of different types} (one truthful day, one lying day).
\begin{recallbox}A truth-teller only says true things; a liar only says false things. After Sunday comes Monday (the week goes round).\end{recallbox}
\step{1}{Table for Chris (L $=$ lies, T $=$ truthful)}
\begin{center}
\setlength{\tabcolsep}{4pt}
\begin{tabular}{|l|c|c|c|c|c|c|c|}
\hline
\rowcolor{tmBlue!10} Day & Mon & Tue & Wed & Thu & Fri & Sat & Sun \\ \hline
Chris today & T & T & T & T & L & L & L \\ \hline
Chris tomorrow & T & T & T & L & L & L & T \\ \hline
Can say it? & -- & -- & -- & \textbf{yes} & -- & -- & \textbf{yes} \\ \hline
\end{tabular}
\end{center}
Chris can say it on \textbf{Thursday} and \textbf{Sunday} (the two ``switching'' days).
\step{2}{Table for Mark}
\begin{center}
\setlength{\tabcolsep}{4pt}
\begin{tabular}{|l|c|c|c|c|c|c|c|}
\hline
\rowcolor{tmBlue!10} Day & Mon & Tue & Wed & Thu & Fri & Sat & Sun \\ \hline
Mark today & T & L & L & L & T & T & T \\ \hline
Mark tomorrow & L & L & L & T & T & T & T \\ \hline
Can say it? & \textbf{yes} & -- & -- & \textbf{yes} & -- & -- & -- \\ \hline
\end{tabular}
\end{center}
Mark can say it on \textbf{Monday} and \textbf{Thursday}.
\step{3}{Find the common day}
The day in both lists: \textbf{Thursday}.
\begin{warnbox}Many students choose Friday, thinking ``Chris lies on Friday''. But on Friday Chris lies and tomorrow (Saturday) he lies too,
so ``Tomorrow I will lie'' is a \emph{true} statement -- and a person who is lying cannot say something true.\end{warnbox}
\verify{Thursday: Chris tells the truth, and tomorrow (Friday) he lies -- the statement is true, so he can say it. Mark lies, and tomorrow
(Friday) he tells the truth -- the statement is false, which fits Mark lying today.}
So the answer is \textbf{(B)}.
\ans{(B) Thursday}
}

% ---------------------------------------------------------------- Problem 9
\begin{probbox}{9}{\relax}
Both scales below are balanced; identical shapes have equal masses.
\begin{center}
\lgaScale{\lgaCir\,\lgaCir\,\lgaCir}{\lgaTri\,\lgaTri}\hspace{1cm}\lgaScale{\lgaSq\,\lgaCir\,\lgaTri}{\lgaSq\,\lgaSq}
\end{center}
Which of the following groups has the same mass as the group \ \lgaCir\,\lgaTri\,\lgaTri\,\lgaTri\ \ (one circle and three triangles)?
\choicesII{\lgaSq\,\lgaSq\,\lgaSq}{\lgaTri\,\lgaTri\,\lgaSq\,\lgaSq}{\lgaCir\,\lgaCir\,\lgaCir\,\lgaCir}{\lgaCir\,\lgaCir\,\lgaCir\,\lgaSq}{\lgaCir\,\lgaCir\,\lgaCir}
\end{probbox}
\probsol{
\textbf{Analysing the problem.} Let the masses of one circle, one triangle and one square be $O$, $T$ and $V$ (in the same unit). The
two scales give two equations. We need the group whose mass equals $O + 3T$.
\begin{recallbox}A balanced scale means the two pans have equal mass. Removing identical objects from \emph{both} pans (subtracting the
same number from both sides) keeps the scale balanced.\end{recallbox}
\step{1}{Write the two equations}
First scale: \ $3O = 2T$. \qquad Second scale: \ $V + O + T = 2V$.
\step{2}{Simplify the equation of the second scale}
Remove one square from each pan (subtract $V$ from both sides): \ $O + T = V$. \emph{One square weighs as much as one circle and one
triangle.}
\step{3}{Introduce an unknown and express every mass in terms of it}
Let $O = x$ \ ($x > 0$). From $2T = 3O = 3x$ we get $T = \dfrac{3x}{2}$. Then
\[ V = O + T = x + \frac{3x}{2} = \frac{5x}{2}. \]
\step{4}{The mass of the given group}
$O + 3T = x + 3\cdot\dfrac{3x}{2} = \dfrac{2x}{2} + \dfrac{9x}{2} = \dfrac{11x}{2}$.
\step{5}{The mass of each option (all written over the denominator $2$ for easy comparison)}
\begin{center}
\renewcommand{\arraystretch}{1.5}
\begin{tabular}{|c|l|c|}
\hline
\rowcolor{tmBlue!10} Option & Mass & Result \\ \hline
(A) & $3V = 3\cdot\dfrac{5x}{2}$ & $\dfrac{15x}{2}$ \\ \hline
(B) & $2T + 2V = 3x + 5x = 8x$ & $\dfrac{16x}{2}$ \\ \hline
(C) & $4O = 4x$ & $\dfrac{8x}{2}$ \\ \hline
(D) & $3O + V = 3x + \dfrac{5x}{2}$ & $\mathbf{\dfrac{11x}{2}}$ \ \checkmark \\ \hline
(E) & $3O = 3x$ & $\dfrac{6x}{2}$ \\ \hline
\end{tabular}
\end{center}
Only option (D) has mass $\dfrac{11x}{2}$.
\step{6}{Another way: substitution only}
$O + 3T = (O + T) + 2T = V + 3O$ \ (replace $O + T$ by $V$ from Step 2, and $2T$ by $3O$ from the first scale). That is exactly three
circles and one square.
\begin{warnbox}The first scale gives $3O = 2T$: we need \emph{more} circles to balance \emph{fewer} triangles, so a circle is
\emph{lighter} than a triangle. Writing $2O = 3T$ by mistake ruins everything from the start.\end{warnbox}
\verify{Take $x = 20$ g: circle $20$ g, triangle $30$ g, square $50$ g. Scale 1: $3\cdot 20 = 60 = 2\cdot 30$; scale 2:
$50 + 20 + 30 = 100 = 2\cdot 50$. Given group: $20 + 3\cdot 30 = 110$ g; option (D): $3\cdot 20 + 50 = 110$ g.}
So the answer is \textbf{(D)}.
\ans{(D) three circles and one square}
}

% ---------------------------------------------------------------- Problem 10
\begin{probbox}{10}{\relax}
Five friends, Chris, Max, Brandon, Kayla and Tanya, go strawberry picking; no two of them pick the same number of strawberries. Chris picks more than
Max. Brandon picks fewer than Kayla. Kayla picks fewer than Max but more than Tanya. Who picks the second most?
\choices{Brandon}{Chris}{Kayla}{Max}{Tanya}
\end{probbox}
\probsol{
\textbf{Analysing the problem.} Five friends, four comparisons of numbers of strawberries. We must find the person in second place (behind
only one other person).
\step{1}{Use symbols (``$>$'' means ``picked more than'')}
$C > M$; \quad $K > B$; \quad $M > K$; \quad $K > T$.
\step{2}{Chain the inequalities}
$C > M > K$, and $K$ is more than both $B$ and $T$.
\begin{figC}
  \node[draw=tmBlue,fill=tmBlue!8,rounded corners=3pt,line width=0.8pt] (H) at (0,0){Chris};
  \node[draw=AnsRed,fill=AnsRed!8,rounded corners=3pt,line width=1.1pt] (L) at (2,0){Max};
  \node[draw=tmBlue,fill=tmBlue!8,rounded corners=3pt,line width=0.8pt] (K) at (4,0){Kayla};
  \node[draw=tmBlue,fill=tmBlue!8,rounded corners=3pt,line width=0.8pt] (B) at (6,0.5){Brandon};
  \node[draw=tmBlue,fill=tmBlue!8,rounded corners=3pt,line width=0.8pt] (T) at (6,-0.5){Tanya};
  \draw[->,Navy,line width=0.8pt] (H)--(L); \draw[->,Navy,line width=0.8pt] (L)--(K);
  \draw[->,Navy,line width=0.8pt] (K)--(B); \draw[->,Navy,line width=0.8pt] (K)--(T);
  \node[font=\scriptsize,text=Slate] at (3,0.75){arrow: ``more than''};
  \draw[<->,dashed,AnsRed!60] (B)--(T) node[midway,right,font=\scriptsize]{?};
\end{figC}
\step{3}{Read off the result}
Chris picks the most (more than Max, who picks more than everyone else). Max is beaten only by Chris and beats Kayla, Brandon, Tanya. So
\textbf{Max picks the second most}. We do not know whether Brandon or Tanya picked more, but that does not affect the question.
\begin{warnbox}The sentence ``Kayla picks fewer than Max but more than Tanya'' contains \emph{two} pieces of information: $M > K$ and $K > T$.
Reading too fast, it is easy to miss one.\end{warnbox}
\verify{For example Chris $30$, Max $25$, Kayla $20$, Brandon $12$, Tanya $10$ strawberries satisfies every condition; Max is second.}
So the answer is \textbf{(D)}.
\ans{(D) Max}
}

% ---------------------------------------------------------------- Problem 11
\begin{probbox}{11}{\relax}
Three friends, Amy, Helen and Tom, count the comics and the science books they read over the summer (no two of them read the same number of
comics, and no two read the same number of science books). Amy read more comics than Helen but fewer science books than Helen. Tom read more
comics than Helen but fewer science books than Amy. Which of the following statements \textbf{must} be true?
\choicesL{Tom read the most comics.}{Amy read the most comics.}{Helen read the fewest comics.}{Amy read the fewest science books.}{Tom read the most science books.}
\end{probbox}
\probsol{
\textbf{Analysing the problem.} There are two completely separate comparisons: \emph{comics} and \emph{science books}. We split the
information into two groups, order each group, then examine each statement.
\begin{recallbox}``Must be true'' $=$ true in \emph{every} situation satisfying the problem. A \emph{single} situation that satisfies the
problem but makes a statement false is enough to rule that statement out.\end{recallbox}
\step{1}{Comics}
``Amy read more comics than Helen'': $A > H$. \quad ``Tom read more comics than Helen'': $T > H$.\\
So Helen read the fewest comics. As for Amy and Tom: \emph{no sentence compares them}.
\step{2}{Science books}
``Amy read fewer science books than Helen'': $H > A$. \quad ``Tom read fewer science books than Amy'': $A > T$.\\
Chained: $H > A > T$. Helen read the most science books, Tom the fewest.
\begin{figC}
  \node[font=\small\bfseries,text=tmOrange!80!black] at (-1,0.6){Comics};
  \node[draw=tmOrange,rounded corners=3pt,line width=0.8pt] (h) at (0.6,0.6){Helen};
  \node[draw=tmOrange,rounded corners=3pt,line width=0.8pt] (a) at (2.4,1.1){Amy};
  \node[draw=tmOrange,rounded corners=3pt,line width=0.8pt] (t) at (2.4,0.1){Tom};
  \draw[->,tmOrange,line width=0.9pt] (h)--(a); \draw[->,tmOrange,line width=0.9pt] (h)--(t);
  \node[font=\scriptsize,text=Slate] at (4.1,0.6){(cannot compare)};
  \node[font=\small\bfseries,text=tmBlue] at (-1,-0.9){Science};
  \node[draw=tmBlue,rounded corners=3pt,line width=0.8pt] (t2) at (0.6,-0.9){Tom};
  \node[draw=tmBlue,rounded corners=3pt,line width=0.8pt] (a2) at (2.4,-0.9){Amy};
  \node[draw=tmBlue,rounded corners=3pt,line width=0.8pt] (h2) at (4.2,-0.9){Helen};
  \draw[->,tmBlue,line width=0.9pt] (t2)--(a2); \draw[->,tmBlue,line width=0.9pt] (a2)--(h2);
  \node[font=\scriptsize,text=Slate] at (6.4,0){arrows point towards ``more''};
\end{figC}
\step{3}{Examine each statement}
\begin{itemize}[leftmargin=1.4em,itemsep=1pt,topsep=1pt]
\item (A) ``Tom read the most comics'': not certain. For example Amy $5$, Tom $3$, Helen $1$ comics satisfies the problem, but Amy read the most.
\item (B) ``Amy read the most comics'': not certain. For example Amy $3$, Tom $5$, Helen $1$ comics satisfies the problem, but Tom read the most.
\item (C) ``Helen read the fewest comics'': \textbf{must be true} (Step 1: Helen read fewer than both Amy and Tom).
\item (D) ``Amy read the fewest science books'': \textbf{false}, since Tom read fewer than Amy.
\item (E) ``Tom read the most science books'': \textbf{false}, Tom read the fewest.
\end{itemize}
\begin{warnbox}Do not mix the two kinds of books into one chain! ``Amy read more comics than Helen'' says nothing about science books. Draw
\emph{two separate diagrams}.\end{warnbox}
So the answer is \textbf{(C)}.
\ans{(C) Helen read the fewest comics}
}

% ---------------------------------------------------------------- Problem 12
\begin{probbox}{12}{\relax}
\begin{minipage}[t]{0.63\linewidth}
Each of the numbers $1, 2, 3, 4, 5, 6, 7, 8$ is placed in a different circle of the triangle shown, so that the sums of the numbers on the
three sides of the triangle (the two slanted sides have $4$ circles each, the base has $3$ circles) are all equal to the same number $S$.
The sum of all the different possible values of $S$ is
\choicesII{$85$}{$99$}{$66$}{$81$}{$67$}
\end{minipage}\hfill
\begin{minipage}[t]{0.33\linewidth}\vspace{0pt}\centering
\lgaTriEight{}{}{}{}{}{}{}{}
\end{minipage}
\end{probbox}
\probsol{
\textbf{Analysing the problem.} Call the number at the top corner $a$, the bottom-left corner $b$, the bottom-right corner $c$; the two
middle numbers of the left side $x_1, x_2$; the two middle numbers of the right side $y_1, y_2$; the middle number of the base $z$. The
three sums:
\[ \text{left: } a + x_1 + x_2 + b = S,\qquad \text{right: } a + y_1 + y_2 + c = S,\qquad \text{base: } b + z + c = S. \]
We must find \emph{every} possible value of $S$ and add them up.
\step{1}{Add the three sides -- how many times is each number counted?}
\begin{recallbox}A corner number belongs to two sides, so it is counted twice; a middle number is counted once.\end{recallbox}
$3S = (1 + 2 + \cdots + 8) + (a + b + c) = 36 + (a + b + c)$, so \ $a + b + c = 3S - 36$.
\step{2}{Pin down $S$}
The corner sum is at least $1 + 2 + 3 = 6$ and at most $6 + 7 + 8 = 21$. Hence $6 \le 3S - 36 \le 21$, i.e.\ $42 \le 3S \le 57$, or
$14 \le S \le 19$. Candidates: $S = 14, 15, 16, 17, 18, 19$.
\step{3}{Use the base to rule some out}
The base has only $3$ numbers: $z = S - (b + c)$. And $b + c = (a + b + c) - a = 3S - 36 - a$, so
\[ z = S - (3S - 36 - a) = 36 - 2S + a. \]
\begin{itemize}[leftmargin=1.4em,itemsep=1pt,topsep=1pt]
\item $S = 14$: $z = 8 + a \ge 9$ -- impossible (the largest number is $8$). \textbf{Ruled out.}
\item $S = 18$: $z = 36 - 36 + a = a$ -- two circles with the same number, not allowed. \textbf{Ruled out.}
\item $S = 15, 16, 17, 19$: not ruled out; we look for actual arrangements.
\end{itemize}
\step{4}{Show arrangements for $S = 15, 16, 17, 19$}
\begin{figC}
  \node at (0,0){\lgaTriEight{$1$}{$4$}{$8$}{$2$}{$3$}{$5$}{$6$}{$7$}}; \node[font=\small] at (0,-1.75){$S = 15$};
  \node at (3.85,0){\lgaTriEight{$1$}{$3$}{$8$}{$4$}{$2$}{$6$}{$7$}{$5$}}; \node[font=\small] at (3.85,-1.75){$S = 16$};
  \node at (7.7,0){\lgaTriEight{$2$}{$3$}{$7$}{$5$}{$1$}{$6$}{$8$}{$4$}}; \node[font=\small] at (7.7,-1.75){$S = 17$};
  \node at (11.55,0){\lgaTriEight{$6$}{$1$}{$5$}{$7$}{$2$}{$3$}{$8$}{$4$}}; \node[font=\small] at (11.55,-1.75){$S = 19$};
\end{figC}
Check, e.g.\ $S = 15$: left $1 + 4 + 8 + 2 = 15$; right $1 + 3 + 5 + 6 = 15$; base $2 + 7 + 6 = 15$. \
$S = 19$: left $6 + 1 + 5 + 7 = 19$; right $6 + 2 + 3 + 8 = 19$; base $7 + 4 + 8 = 19$.
\step{5}{Add the possible values}
$15 + 16 + 17 + 19 = 67$.
\begin{warnbox}If you only pin down $14 \le S \le 19$ and add all six values, you get $99$ -- option (B), a ``trap''. Pinning down only gives
\emph{candidates}; each candidate must be checked: either show an arrangement or prove that none exists.\end{warnbox}
\begin{parentbox}Let your child cut out $8$ slips of paper numbered $1$ to $8$ and try placing them on the figure. After a few attempts, ask:
``Why can we never manage $S = 18$?'' -- that is when your child sees the value of the argument in Step 3.\end{parentbox}
So the answer is \textbf{(E)}.
\ans{(E) $67$}
}

% ---------------------------------------------------------------- Problem 13
\begin{probbox}{13}{\relax}
\begin{minipage}[t]{0.62\linewidth}
The $4\times 4$ grid shown is to be filled so that each row and each column contains the digits $1, 2, 3, 4$. The grid is divided into four
$2\times 2$ squares (thick borders), and each of these squares must also contain the digits $1, 2, 3, 4$. Which digit goes in cell $P$?
\choicesL{$1$}{$2$}{$3$}{$4$}{The digit cannot be determined}
\end{minipage}\hfill
\begin{minipage}[t]{0.33\linewidth}\vspace{4pt}\centering
\lgaFour{1/1/$1$,3/1/$3$,2/2/$2$,2/3/$P$,4/4/$4$}
\end{minipage}
\end{probbox}
\probsol{
\textbf{Analysing the problem.} This is a $4\times 4$ sudoku. Cell $P$ is in row $3$, column $2$. We reason step by step; every digit we
write must have a reason.
\begin{recallbox}Each cell belongs to one row, one column and one $2\times 2$ box; its digit must differ from every digit already in those three regions.\end{recallbox}
\step{1}{Top-left box}
This box already has $1$ (row $1$, column $1$) and $2$ (row $2$, column $2$); it still needs $3$ and $4$ in two cells: row $1$ column $2$
and row $2$ column $1$. Row $1$ already has a $3$, so row $1$ column $2$ is \textbf{$4$}, and row $2$ column $1$ is \textbf{$3$}.
\step{2}{Column $2$ leaves only two possibilities for $P$}
Column $2$ already has $4$ and $2$, so $P \in \{1, 3\}$.
\step{3}{Rows $1$ and $2$ on the right}
Row $1$: $1, 4, 3, \_$ $\Rightarrow$ row $1$ column $4$ is $2$. \
Row $2$: $3, 2, \_, \_$ still needs $\{1, 4\}$. Column $4$ already has a $4$ (in row $4$), so row $2$ column $4$ is $1$ and row $2$ column $3$ is $4$.
\step{4}{Column $4$ and row $3$}
Column $4$: $2, 1, \_, 4$ $\Rightarrow$ row $3$ column $4$ is \textbf{$3$}. So row $3$ already contains a $3$, hence $P \ne 3$. With Step 2:
$\mathbf{P = 1}$.
\begin{figC}
  \node at (0,0){\lgaFour{1/1/$1$,2/1/\textcolor{tmGreen!60!black}{$4$},3/1/$3$,4/1/\textcolor{tmGreen!60!black}{$2$},1/2/\textcolor{tmGreen!60!black}{$3$},2/2/$2$,3/2/\textcolor{tmGreen!60!black}{$4$},4/2/\textcolor{tmGreen!60!black}{$1$},2/3/\textcolor{AnsRed}{$\mathbf{1}$},4/3/\textcolor{tmGreen!60!black}{$3$},4/4/$4$}};
  \node at (4.2,0){\lgaFour{1/1/$1$,2/1/$4$,3/1/$3$,4/1/$2$,1/2/$3$,2/2/$2$,3/2/$4$,4/2/$1$,1/3/$4$,2/3/\textcolor{AnsRed}{$\mathbf{1}$},3/3/$2$,4/3/$3$,1/4/$2$,2/4/$3$,3/4/$1$,4/4/$4$}};
  \node[font=\scriptsize,text=Slate] at (0,-1.75){cells just deduced};
  \node[font=\scriptsize,text=Slate] at (4.2,-1.75){completed grid};
\end{figC}
\step{5}{Complete the grid to be sure there is no contradiction}
Column $1$ still needs $\{2, 4\}$; row $4$ has a $4$, so row $4$ column $1$ is $2$ and row $3$ column $1$ is $4$. Row $3$: $4, 1, \_, 3$ $\Rightarrow$ $2$.
Row $4$: $2, \_, \_, 4$ needs $\{1, 3\}$; column $2$ still needs a $3$, so row $4$ column $2$ is $3$ and column $3$ is $1$.
\verify{Every row, column and $2\times 2$ box of the completed grid contains all of $1, 2, 3, 4$. The grid can be filled in only one way, so $P$ is determined.}
\begin{warnbox}Option (E) is a trap for those who look only at column $2$ (seeing two possibilities, $1$ and $3$) and stop there. Keep going in other regions!\end{warnbox}
So the answer is \textbf{(A)}.
\ans{(A) $1$}
}
